\documentclass[12pt,a4paper]{article}
\usepackage[margin=1in]{geometry}
\usepackage{amsmath,amssymb,amsthm}
\usepackage{array}
\usepackage{booktabs}
\usepackage{hyperref}
\usepackage{enumitem}
\usepackage{algorithm}
\usepackage{algpseudocode}
\usepackage{graphicx}
\usepackage[numbers]{natbib}

\hypersetup{
  colorlinks=true,
  linkcolor=blue,
  citecolor=blue,
  urlcolor=blue
}

\newtheorem{theorem}{Theorem}
\newtheorem{lemma}[theorem]{Lemma}
\newtheorem{proposition}[theorem]{Proposition}
\newtheorem{definition}{Definition}

\title{Complexity Analysis and Optimization of Cloud-Based\\Duplicate Folder Reconciliation in Aethel}
\author{%
  Aethel Project\\
  \texttt{https://github.com/aethel}
}
\date{April 2026}

\begin{document}
\maketitle

\begin{abstract}
Cloud storage synchronization systems routinely encounter duplicate directory structures arising from concurrent edits, interrupted uploads, or multi-device conflicts. The \textsc{Aethel} system implements a deduplication pipeline (\texttt{dedupe-folders~--execute}) that detects and merges structurally equivalent folder subtrees on Google Drive. In this paper, we present a formal complexity analysis of the current reconciliation algorithm, identify the dominant cost factors---particularly the amplification effect of remote API latency on per-item operations---and propose a five-part optimization plan that reduces the asymptotic operation count from $O(P \cdot N + K)$ to $O(N + K)$ while significantly improving wall-clock performance through structural batching, depth-first traversal ordering, and controlled I/O parallelism. We validate our analysis against a representative case study involving deeply nested, high-fan-out subtrees such as Python virtual environments and IDE build artifacts.
\end{abstract}

\section{Introduction}

Deduplication of hierarchical namespaces is a well-studied problem in file systems and distributed storage~\cite{quinlan2002venti, zhu2008avoiding}. However, when the target storage layer is a rate-limited cloud API rather than a local file system, the cost model shifts fundamentally: the dominant expense is no longer CPU-bound comparison but network round-trip latency per metadata operation.

\textsc{Aethel} is a synchronization tool that manages a local--remote mirror against Google Drive. During routine operation, structural inconsistencies---specifically, sibling directories sharing identical names within the same parent---accumulate due to concurrent modifications across devices or interrupted synchronization sessions. The \texttt{dedupe-folders~--execute} command performs a multi-pass reconciliation that detects duplicate folder groups, selects a canonical representative, and migrates children from redundant folders into the canonical target.

The primary implementation resides in \texttt{src/core/drive-api.js}, and the algorithm can be decomposed into five principal stages:

\begin{enumerate}[nosep]
  \item \textbf{Enumeration}: A full listing of all remote items under the synchronization root via \texttt{fetchAllItems()}.
  \item \textbf{Detection}: Construction of duplicate folder groups through \texttt{buildDuplicateFolderGroups()}, which identifies sibling folders sharing the same name.
  \item \textbf{Selection}: For each duplicate group, election of a \emph{canonical folder}; remaining members are designated as \emph{loser folders}.
  \item \textbf{Merge}: Recursive migration of children from each loser folder into the canonical folder via \texttt{mergeFolderIntoCanonical()}.
  \item \textbf{Convergence}: Re-scanning the remote tree to detect newly exposed duplicates; if any exist, the process iterates from Stage~1.
\end{enumerate}

The remainder of this paper is organized as follows. Section~\ref{sec:notation} establishes notation. Section~\ref{sec:analysis} provides a formal complexity analysis under both computational and I/O cost models. Section~\ref{sec:bottleneck} examines why wall-clock time diverges substantially from the theoretical bound. Section~\ref{sec:casestudy} presents a concrete case study. Section~\ref{sec:optimization} proposes a detailed optimization plan. Section~\ref{sec:conclusion} summarizes our findings.


\section{Notation and Definitions}
\label{sec:notation}

\begin{definition}[Remote Item Set]
Let $\mathcal{I}$ denote the set of all remote items (files and folders) under the synchronization root, with $N = |\mathcal{I}|$.
\end{definition}

\begin{definition}[Duplicate Affected Items]
Let $K$ denote the total number of child items across all duplicate folder subtrees that require inspection, comparison, migration, or deletion during the reconciliation process.
\end{definition}

\begin{definition}[Convergence Depth]
Let $P$ denote the number of full passes required for the reconciliation process to reach a fixed point---i.e., no further duplicate sibling folders exist. When merging an upper-level duplicate group exposes previously hidden lower-level duplicates, $P > 1$.
\end{definition}

\begin{center}
\begin{tabular}{c p{10cm}}
\toprule
\textbf{Symbol} & \textbf{Definition} \\
\midrule
$N$ & Total number of remote items under the synchronization root \\
$K$ & Aggregate count of child items within duplicate subtrees requiring processing \\
$P$ & Number of convergence passes until no duplicates remain \\
$\lambda$ & Mean round-trip latency per Google Drive API call \\
$B$ & Page size of \texttt{files.list} responses (typically 1000) \\
\bottomrule
\end{tabular}
\end{center}


\section{Complexity Analysis}
\label{sec:analysis}

\subsection{Computational Complexity}

Considering only local computation (in-memory comparisons, hash lookups, tree traversals), the total cost of the reconciliation algorithm is:

\begin{equation}
T(N, K, P) = O(P \cdot N + K)
\label{eq:time}
\end{equation}

The $O(P \cdot N)$ term arises from re-executing a complete remote enumeration and duplicate group construction in each pass. The $O(K)$ term accounts for per-child operations within duplicate subtrees: name lookup against the canonical folder's children, parent pointer reassignment (move), or deletion of exact duplicates.

\subsection{I/O Operation Complexity}

The computational complexity in Equation~\eqref{eq:time} substantially understates the true cost because the bottleneck is not local CPU cycles but remote API invocations, each incurring latency $\lambda$. A more faithful cost model counts discrete API operations:

\begin{equation}
A(N, K, P) \approx P \cdot \left\lceil \frac{N}{B} \right\rceil + C_{\text{children}} + C_{\text{lookup}} + C_{\text{write}}
\label{eq:api}
\end{equation}

where:
\begin{itemize}[nosep]
  \item $P \cdot \lceil N/B \rceil$ counts paginated \texttt{files.list} calls across all passes,
  \item $C_{\text{children}}$ counts \texttt{listDirectChildren()} invocations---at least two per loser folder (one to enumerate children, one to verify the folder is empty post-merge),
  \item $C_{\text{lookup}}$ counts per-child name-collision queries against the canonical parent,
  \item $C_{\text{write}}$ counts \texttt{files.update} (move) and \texttt{files.trash} (delete duplicate) mutations.
\end{itemize}

\begin{proposition}
The expected wall-clock time is lower-bounded by:
\[
W \geq \lambda \cdot A(N, K, P)
\]
Under sequential execution with no request pipelining, this bound is tight up to constant factors.
\end{proposition}

Since $\lambda$ is typically 100--500\,ms for Google Drive API calls, even moderate values of $A$ translate to substantial elapsed time.


\section{Analysis of Wall-Clock Divergence}
\label{sec:bottleneck}

The observed wall-clock time exceeds what the asymptotic bound $O(P \cdot N + K)$ might suggest for an equivalent in-memory algorithm by several orders of magnitude. We identify four principal contributing factors:

\paragraph{1. Network-Dominated Execution Profile.}
Empirical observation shows near-zero CPU utilization during reconciliation, with continuous output indicating active waiting on API responses. The process is fundamentally I/O-bound, and its throughput is governed by $\lambda$ rather than instruction count.

\paragraph{2. Redundant Full-Tree Enumeration.}
Each convergence pass re-fetches the entire remote tree from scratch. For large $N$, the enumeration phase alone is expensive: at $B = 1000$ items per page and $\lambda = 200$\,ms, a tree with $N = 50{,}000$ items requires $\sim$10\,s per pass solely for listing.

\paragraph{3. Per-Item Query Granularity.}
The current implementation does not batch children of the same parent into a single in-memory index. Instead, each child item may trigger an independent API call to check for name collisions in the target folder, yielding $O(K)$ sequential API round-trips.

\paragraph{4. Multi-Pass Convergence.}
The algorithm's top-down processing order means that merging an upper-level duplicate group may create new sibling collisions at deeper levels, necessitating additional passes. This coupling between levels inflates $P$ beyond the minimum achievable.


\section{Case Study}
\label{sec:casestudy}

The motivating case involves a synchronization root containing several high-fan-out, deeply nested subtrees:

\begin{itemize}[nosep]
  \item \textbf{Python virtual environments} (\texttt{venv/site-packages/}): These directories contain thousands of small files across hundreds of packages, inflating $K$ disproportionately.
  \item \textbf{Build artifacts}: Compiler output trees, IDE metadata directories (e.g., \texttt{.idea/}, \texttt{.vscode/}), and Xcode project structures introduce deeply nested duplicate paths.
  \item \textbf{Course and project artifacts}: Educational material repositories with duplicated assignment structures.
\end{itemize}

The distinguishing characteristic of this workload is not large individual file sizes but an exceptionally high item count composed of small files. Since API latency per operation is approximately constant regardless of file size, the total cost scales with the \emph{number} of items rather than their aggregate size. This renders the problem a \emph{metadata-intensive reconciliation} task rather than a data-transfer-dominated synchronization problem.

Furthermore, the depth of duplication across these subtrees causes $P$ to exceed unity by a significant margin, as each merge pass exposes previously shadowed duplicates deeper in the hierarchy.


\section{Proposed Optimization Plan}
\label{sec:optimization}

We propose a five-part optimization strategy targeting both algorithmic complexity and I/O efficiency. The strategies are ordered by expected impact and should be implemented incrementally.

\subsection{Strategy 1: Single-Pass Tree Materialization}

\paragraph{Current behavior.} Each convergence pass invokes \texttt{fetchAllItems()} to retrieve the full remote tree, constructing the duplicate group index from scratch.

\paragraph{Proposed change.} Perform exactly one full enumeration at initialization. Construct an in-memory \emph{folder graph} $G = (V, E)$ where vertices represent folders and edges represent parent--child relationships, along with a \emph{sibling index} $\sigma: (\text{parentId}, \text{name}) \to \mathcal{P}(\text{itemId})$ mapping each (parent, name) pair to the set of items sharing that key.

After each merge operation, update $G$ and $\sigma$ incrementally rather than re-fetching the entire tree:

\begin{algorithmic}[1]
\Function{IncrementalUpdate}{$G$, $\sigma$, movedItem, newParent}
  \State Remove movedItem from old parent's child set in $G$
  \State Update $\sigma$ to remove movedItem from old (parent, name) entry
  \State Add movedItem to newParent's child set in $G$
  \State Update $\sigma$ to insert movedItem at new (parent, name) entry
\EndFunction
\end{algorithmic}

\paragraph{Complexity reduction.} This eliminates the $P$ multiplier on the enumeration cost, reducing the listing component from $O(P \cdot N)$ to $O(N)$.

\subsection{Strategy 2: Depth-First (Deepest-First) Merge Ordering}

\paragraph{Current behavior.} Duplicate groups are processed in discovery order, which is typically breadth-first or unordered. Upper-level merges may expose new duplicates at lower levels, requiring additional passes.

\paragraph{Proposed change.} Sort duplicate groups by depth in descending order (deepest first). Process all groups at depth $d$ before any group at depth $d - 1$.

\begin{proposition}
Under deepest-first ordering with the single-pass tree materialization of Strategy~1, the convergence depth satisfies $P = 1$---i.e., a single logical pass suffices.
\end{proposition}

\begin{proof}[Proof sketch]
Merging at depth $d$ cannot create new duplicates at depth $d' > d$, since no items are moved to deeper levels. By processing all groups at the maximum depth first and proceeding upward, each merge at depth $d$ finds the subtrees below already reconciled. Thus no second pass is required. \qed
\end{proof}

\paragraph{Complexity reduction.} Reduces $P$ to 1, yielding an overall bound of $O(N + K)$.

\subsection{Strategy 3: Sibling Child Cache}

\paragraph{Current behavior.} When migrating a child from a loser folder to the canonical folder, the implementation queries the canonical folder's children to detect name collisions on a per-item basis.

\paragraph{Proposed change.} Before processing any loser folder under a given canonical parent, fetch the canonical parent's full child list once and index it as a hash map $\mu: \text{name} \to \text{itemId}$. All subsequent collision checks become $O(1)$ local lookups.

\paragraph{Complexity reduction.} Eliminates $C_{\text{lookup}}$ from the API operation count entirely, replacing it with a single \texttt{listDirectChildren()} call per canonical folder.

\subsection{Strategy 4: Controlled I/O Parallelism}

\paragraph{Current behavior.} All API operations are executed sequentially.

\paragraph{Proposed change.} Introduce bounded concurrency for independent operations. Two levels of parallelism are safe to exploit:

\begin{enumerate}[nosep]
  \item \textbf{Inter-group parallelism}: Duplicate groups whose subtrees do not share any ancestor--descendant relationship can be processed concurrently.
  \item \textbf{Intra-group parallelism}: Within a single merge operation, independent child migrations (where no two children share the same target name) can be issued concurrently.
\end{enumerate}

Concurrency should be bounded by a semaphore with limit $c$ (e.g., $c = 5$--$10$) to remain within Google Drive API rate limits (approximately 10~queries/second/user for mutation operations).

\paragraph{Complexity reduction.} Does not change the asymptotic operation count, but reduces wall-clock time by a factor approaching $\min(c, A/\lambda)$ in the I/O-bound regime.

\subsection{Strategy 5: Synchronization Boundary Exclusion}

\paragraph{Current behavior.} All items under the synchronization root are enumerated and subject to deduplication, including directories whose contents are reproducible and should not be synchronized.

\paragraph{Proposed change.} Introduce a configurable exclusion list (analogous to \texttt{.gitignore}) that prunes specified subtrees from both enumeration and deduplication:

\begin{verbatim}
  # Suggested default exclusions
  venv/
  node_modules/
  __pycache__/
  .idea/
  .vscode/
  build/
  dist/
  *.xcodeproj/
\end{verbatim}

\paragraph{Complexity reduction.} Reduces both $N$ and $K$ by excluding high-fan-out, low-value subtrees. In the motivating case study, exclusion of \texttt{venv/} alone would eliminate the majority of affected items.


\subsection{Projected Improvement}

With all five strategies implemented, the optimized reconciliation achieves:

\begin{equation}
T_{\text{opt}}(N', K') = O(N' + K')
\label{eq:optimized}
\end{equation}

where $N' \leq N$ and $K' \leq K$ reflect the reduced item counts after boundary exclusion. The convergence depth is reduced to $P = 1$, redundant API calls are eliminated through caching, and wall-clock time is further compressed by a parallelism factor of up to $c$.

\begin{center}
\begin{tabular}{l c c}
\toprule
\textbf{Optimization} & \textbf{Target Cost Factor} & \textbf{Expected Reduction} \\
\midrule
Single-pass materialization & $P \cdot \lceil N/B \rceil$ & Eliminates factor $P$ \\
Deepest-first ordering & $P$ & Guarantees $P = 1$ \\
Sibling child cache & $C_{\text{lookup}}$ & Eliminates entirely \\
Controlled parallelism & Wall-clock time & $\sim c\times$ speedup \\
Boundary exclusion & $N$, $K$ & Problem-dependent; often $>50\%$ \\
\bottomrule
\end{tabular}
\end{center}


\section{Conclusion}
\label{sec:conclusion}

The current implementation of \texttt{dedupe-folders~--execute} in \textsc{Aethel} exhibits a time complexity of $O(P \cdot N + K)$, where the dominant cost is not local computation but remote API latency amplified across thousands of sequential metadata operations. The performance degradation observed in practice stems from four compounding factors: redundant full-tree enumeration across convergence passes, per-item remote queries without local caching, top-down processing order that inflates convergence depth, and the absence of I/O parallelism.

The proposed optimization plan addresses each factor systematically. Single-pass tree materialization and deepest-first merge ordering together reduce the asymptotic bound to $O(N + K)$ with guaranteed single-pass convergence. Sibling caching eliminates redundant lookup operations. Controlled parallelism exploits the I/O-bound nature of the workload to improve throughput by a constant factor. Finally, synchronization boundary exclusion reduces the problem size at its source by filtering reproducible, high-fan-out subtrees before enumeration.

Collectively, these optimizations transform the reconciliation from a multi-pass, sequentially I/O-bound process into a single-pass, partially parallel procedure operating over a reduced item set---addressing both the theoretical complexity and the practical wall-clock performance that ultimately determines user experience.


\bibliographystyle{plainnat}
\begin{thebibliography}{9}

\bibitem[Quinlan and Dorward(2002)]{quinlan2002venti}
S.~Quinlan and S.~Dorward.
\newblock Venti: A new approach to archival storage.
\newblock In \emph{Proceedings of the USENIX Conference on File and Storage Technologies (FAST)}, pages 89--101, 2002.

\bibitem[Zhu et~al.(2008)]{zhu2008avoiding}
B.~Zhu, K.~Li, and H.~Patterson.
\newblock Avoiding the disk bottleneck in the data domain deduplication file system.
\newblock In \emph{Proceedings of the USENIX Conference on File and Storage Technologies (FAST)}, pages 269--282, 2008.

\end{thebibliography}

\end{document}
